% ID: MAT-ALG-INT-005
% Titolo: Numeri interi e proprietà delle potenze 5
% Disciplina: Matematica
% Area: Algebra
% Argomento: Numeri interi
% Classe: Prima secondaria di secondo grado
% Difficolta: 3
% Tipo: problema_numerico
% Tag: matematica, algebra, numeri_interi
% Fonte: verifica_interi_anonimizzata
% Provenienza: materiale fornito dal docente; testo anonimizzato; soluzione redatta per Exergo

\begin{esercizio}
\[
    (-3)^3\cdot[(-9^2)^5:(-3^{3^2}):(-3^3)^2]:(-9^{2^2})
    \cdot
    \left\{[-(66^3:11^3)^2]^3:[(-6^{2^2})\cdot(6^5)^2\cdot(-36)]\right\}
    \]
\end{esercizio}

\begin{soluzione}
La prima parentesi quadra vale $(-3^{20})/[(-3^9)3^6]=3^5$.
Dividendo per $-9^4=-3^8$ e moltiplicando per $-3^3$ si ottiene 1.
Nella parentesi graffa il numeratore è $-6^{18}$ e il denominatore $6^{16}$.
Il risultato complessivo è $-36$. La potenza precede il segno meno.
\end{soluzione}
